\chapter{Полная графовая конкуренция конденсата когерентности}
\label{ch:v7-edge-coherence-full-graph-competition-gate}

\section{Возврат к точному меню рёбер}

Полевой суперсвязностный гейт доказал существование вспомогательного
носителя для матрицы $B\in M_{2\times3}(\mathbb C)$. Теперь необходимо
проверить более сильное физическое утверждение: совпадает ли ранга-один
вакуум $B$ с шестью новыми рёбрами ранее выбранного вектороподобного
ремонта.

Полный строгий граф содержит три старых и одиннадцать новых рёбер. Поле
$B$ расположено на полном двудольном прямоугольнике
\begin{equation}
 \{L_L,Y_L\}\times\{e_R,X_R,Y_R\}.
 \label{eq:v7-full-graph-coherence-biclique}
\end{equation}
В порядке строк $(L_L,Y_L)$ и столбцов $(e_R,X_R,Y_R)$ целевое меню внутри
этого прямоугольника имеет маску
\begin{equation}
 M_{\rm target}^{B}
 =\begin{pmatrix}1&0&1\\0&0&1\end{pmatrix}.
 \label{eq:v7-full-graph-target-mask}
\end{equation}
Здесь $(L_L,e_R)$ --- старое ребро, а $(L_L,Y_R)$ и $(Y_L,Y_R)$ --- два
новых целевых ребра. Три остальных новых ребра прямоугольника относятся к
нежелательным.

Ещё шесть новых рёбер лежат вне $B$:
\begin{equation}
 \{Q_LY_R,\ X_LX_R,\ X_LY_R,\ X_Ld_R,\ X_Le_R,\ X_Lu_R\}.
 \label{eq:v7-full-graph-outside-edges}
\end{equation}
Из них четыре входят в целевой ремонт, а два являются нежелательными.

\section{Точный ранговый конфликт}

Общий представитель целевой опоры внутри $B$ имеет вид
\begin{equation}
 B_{\rm target}
 =\begin{pmatrix}a&0&b\\0&0&c\end{pmatrix},
 \qquad abc\ne0.
 \label{eq:v7-full-graph-target-field}
\end{equation}
Его минор по столбцам $e_R,Y_R$ равен
\begin{equation}
 \det B_{\{e,Y\}}=ac\ne0,
 \qquad \rank B_{\rm target}=2.
 \label{eq:v7-full-graph-target-rank-two}
\end{equation}
Но нулевое множество спектрального потенциала когерентности требует
\begin{equation}
 \Lambda^2B=0,
 \qquad \rank B=1.
 \label{eq:v7-full-graph-coherence-rank-one}
\end{equation}
Следовательно, точная целевая опора и вакуум когерентности не пересекаются.

То же препятствие видно без выбора амплитуд. Опора ненулевой матрицы ранга
один есть декартово произведение опор двух векторов:
\begin{equation}
 \operatorname{supp}(uv^\dagger)
 =\operatorname{supp}(u)\times\operatorname{supp}(v).
 \label{eq:v7-full-graph-rank-one-rectangle}
\end{equation}
Среди всех $21$ ненулевых прямоугольных опор размера $2\times3$ целевой
рисунок отсутствует. Если сохранить все три требуемых ребра
$(L_L,e_R)$, $(L_L,Y_R)$ и $(Y_L,Y_R)$, прямоугольность принудительно
включает нежелательное ребро $(Y_L,e_R)$.

\section{Отсутствие целевой стационарной точки}

Положим $x=|a|^2$, $y=|b|^2$, $z=|c|^2$. На целевой страте действие
ограничивается к
\begin{equation}
 \mathcal S_B
 =(x+y+z-3)^2+\frac53xz.
 \label{eq:v7-full-graph-target-restricted-action}
\end{equation}
Уравнение по $y$ требует $x+y+z-3=0$. После этого уравнение по $x$
требует $z=0$. Поэтому внутри открытой страты $xyz\ne0$ стационарной
точки нет. Поток уводит целевую конфигурацию на границу, где исчезает как
минимум одно необходимое ребро.

Для нормированного представителя $a=b=c=1$ имеем
\begin{equation}
 \Tr BB^\dagger=3,qquad
 \det(BB^\dagger)=1,qquad
 \mathcal S_B=\frac53>0.
 \label{eq:v7-full-graph-target-energy}
\end{equation}

\section{Шесть внешних спектаторов}

Допущенная суперсвязность $\mathcal D_B$ зависит только от шести
компонент прямоугольника $B$. Если остальные шесть новых комплексных рёбер
обозначить $z\in\mathbb C^6$, то минимальное расширение имеет вид
\begin{equation}
 \mathcal S_{\rm ext}(B,z)=\mathcal S_B(B),
 \qquad \frac{\partial\mathcal S_{\rm ext}}{\partial z}=0.
 \label{eq:v7-full-graph-spectator-action}
\end{equation}
В одном поколении это добавляет двенадцать вещественных нулевых направлений;
для шести независимых матриц $3\times3$ --- сто восемь. Полный
однопоколенный гессиан в ранга-один минимуме имеет сигнатуру
\begin{equation}
 (n_-,n_0,n_+)=(0,19,5).
 \label{eq:v7-full-graph-extended-hessian}
\end{equation}
Среди спектаторов остаются четыре нужных и два нежелательных новых ребра.
Следовательно, текущее действие не выбирает ни их наличие, ни их отсутствие.

\begin{theorem}[Запрет когерентности ранга один как селектора целевого графа]
Положительный полевой суперсвязностный носитель цепи $1\to6\to3$ не
выбирает шестирёберное вектороподобное расширение. Целевая опора внутри
его прямоугольника имеет ранг два и не является стационарной точкой
ранга-один потенциала. Любой ранга-один представитель, сохраняющий три
нужных внутренних ребра, включает хотя бы одно нежелательное. Остальные
шесть новых рёбер отсутствуют в действии и остаются плоскими.
\end{theorem}

Математический успех предыдущего гейта сохраняется: он даёт корректный
механизм ранга-один когерентности. Закрывается только его более сильное
чтение как селектора конкретного физического меню рёбер.

\begin{nogocriterion}[Запрет ручного исправления опоры]
Нельзя занулить три нежелательных компоненты $B$, добавить положительные
массы двум внешним нежелательным рёбрам и отдельно сохранить четыре внешних
целевых ребра. Такой набор операций и есть искомый селектор, записанный в
действие вручную. Новый родитель обязан видеть весь строгий колчан и
выводить целевой цикл одной кривизностью или одной системой relations.
\end{nogocriterion}

\begin{protocol}
\begin{enumerate}
 \item новых разрешённых рёбер: \textbf{$11$};
 \item новых рёбер внутри $B$: \textbf{$5$};
 \item целевых новых рёбер внутри $B$: \textbf{$2$};
 \item нежелательных новых рёбер внутри $B$: \textbf{$3$};
 \item новых рёбер вне $B$: \textbf{$6$};
 \item целевых новых рёбер вне $B$: \textbf{$4$};
 \item нежелательных новых рёбер вне $B$: \textbf{$2$};
 \item ранг целевой внутренней опоры: \textbf{$2$};
 \item ранга-один целевая опора существует: \textbf{нет};
 \item целевая внутренняя стационарная точка: \textbf{нет};
 \item принудительное нежелательное ребро: \textbf{$Y_L-e_R$};
 \item внешних вещественных нулевых мод: \textbf{$12$ на поколение};
 \item вспомогательный суперсвязностный носитель: \textbf{сохранён};
 \item селектор целевого шестирёберного ремонта: \textbf{закрыт};
 \item следующий гейт: \textbf{полноколчанный родитель кривизны цикла}.
\end{enumerate}
\end{protocol}

Машинный сертификат сохранён в
\path{s2t_v7_edge_coherence_full_graph_competition_gate_results.json}.

\begin{thebibliography}{9}
\bibitem{v7-full-graph-determinantal}
F.~J.~Kir\'aly, L.~Theran, R.~Tomioka,
\emph{The algebraic combinatorial approach for low-rank matrix completion},
Journal of Machine Learning Research 16 (2015), 1391--1436;
arXiv:1211.4116.
\bibitem{v7-full-graph-quiver-morse}
M.~Harada, G.~Wilkin,
\emph{Morse theory of the moment map for representations of quivers},
Geometriae Dedicata 150 (2011), 307--353; arXiv:0807.4734.
\end{thebibliography}